\section{Convergence dynamics from finite energy}
\label{sec:dynamics}

We retain the Cayley graph $G$, the compactification $X=\Gamma\sqcup B$,
and its metric $\rho$ from the preceding section. Write
$\lambda(e)=\rho(e^-,e^+)$ for the length of an edge. Thus
$\lambda\in\ell^2(E(G))$. We will show that translations collapse the
regular coordinate cuts. Two such cuts then detect the attracting and
repelling points of every escaping sequence of group elements.

First observe that, for every fixed $s\in\Gamma$,
\begin{equation}
 \rho(u,us)\longrightarrow0\qquad(u\longrightarrow\infty\text{ in }\Gamma).
 \label{eq:dyn-right-move}
\end{equation}
Here and below, tending to infinity in a discrete group means leaving
every finite subset. Indeed, translate a fixed finite edge path from $1$
to $s$ by $u$. Each translated edge escapes every finite edge set, and an
$\ell^2$ function tends to zero outside finite sets. Summing the finitely
many edge lengths proves \eqref{eq:dyn-right-move}. The convergence is
uniform over any fixed finite set of choices of $s$.

\subsection{Lengths of translated cuts}

Recall that $S(F,t)$ consists of the edges whose endpoint values of $F$
lie on opposite sides of $t$. We always exclude levels attained at
vertices. Call two edges near if they belong to the boundary of a common
cell, or are equal. Enlarging this relation by a fixed graph distance
would not affect the following argument. Its incidence is uniformly
bounded in either variable.

\begin{lemma}[Translated cut lengths]\label{dyn:length}
For every base coordinate $F$ and $g\in\Gamma$,
\begin{equation}
 \int_0^1\sum_{e\in S(F,t)}\sum_{e'\text{ near }e}
       \lambda(ge')\,dt
 =\sum_e |dF(e)|\sum_{e'\text{ near }e}\lambda(ge').
 \label{eq:dyn-length}
\end{equation}
This quantity is finite and tends to zero as $g\to\infty$.
There is a constant $C$, independent of $g,F,t$, such that for every
dually connected set $D$ of edges,
\begin{equation}
 \operatorname{diam}_{\rho}\{gv:v\text{ is an endpoint of an edge of }D\}
 \le C\sum_{e\in D}\sum_{e'\text{ near }e}\lambda(ge').
 \label{eq:dyn-dual-diameter}
\end{equation}
The sum of these diameter bounds over distinct dual components is
bounded by the same expression with their union in place of $D$.
\end{lemma}

\begin{proof}
For each edge, the set of straddled levels has Lebesgue measure
$|dF(e)|$. Tonelli's theorem gives \eqref{eq:dyn-length}. Bounded
incidence implies that
\[
 a(e')=\sum_{e:\ e'\text{ near }e}|dF(e)|
\]
belongs to $\ell^2(E(G))$. The right side of
\eqref{eq:dyn-length} is $\sum_{e'}a(e')\lambda(ge')$, so it is finite
by Cauchy--Schwarz. Translates of a fixed finite edge set eventually
avoid any other fixed finite edge set. Approximate both $a$ and
$\lambda$ in $\ell^2$ by finitely supported functions and apply
Cauchy--Schwarz to the errors. Their translated pairing therefore
tends to zero.

To prove \eqref{eq:dyn-dual-diameter}, join any two edges of $D$ by a
simple path in the dual adjacency graph. For each consecutive pair,
traverse the boundary of a cell containing both edges. This joins
their endpoints using edges near $D$. The simple dual path uses each
of its edges at most once; bounded cell size and incidence give a
uniform bound on the multiplicity with which any near edge is used.
The triangle inequality proves the estimate. Sum it over the
components to obtain the last assertion.
\end{proof}

\begin{lemma}[Collapse of translated cuts]\label{dyn:walls}
Suppose $g_n\to a\in B$. Given finitely many coordinates from the
enlarged coordinate family, there is a subsequence such that, for
almost every level of each coordinate, all endpoints of the
$g_n$-translated straddling set tend uniformly to $a$.
If also $g_n^{-1}\to b\in B$, the subsequence can be chosen so that the
corresponding inverse-translated straddling sets tend uniformly to $b$.
\end{lemma}

\begin{proof}
For a finite list of base functions, Lemma~\ref{dyn:length} permits a
subsequence along which the sum over $n$ of all the integrals in
\eqref{eq:dyn-length} is finite. For almost every level, the
corresponding nonnegative length sums then tend to zero. Each regular
base cut has finitely many dual components. Choose one endpoint in
each nonempty component. The translates of these finitely many fixed
vertices tend to $a$ by \eqref{eq:dyn-right-move}. The diameter estimate
in Lemma~\ref{dyn:length} therefore makes every endpoint of the
translated cut tend uniformly to $a$.

For an added coordinate associated with a box, its straddling set is
contained in a finite union of base straddling sets at levels of the
form $a_i+t$ and $b_i-t$. Include those base coordinates in the list.
Translations and reflections of the level parameter preserve null
sets, so almost every $t$ has the required estimates for all the
covering cuts. Levels outside $[0,1]$ give empty cuts. This proves the
assertion for added coordinates. To obtain both directions, choose
the summable subsequence simultaneously for $g_n$ and $g_n^{-1}$.
\end{proof}

\subsection{The convergence action}

\begin{proposition}\label{dyn:convergence}
If $g_n\to a\in B$ and $g_n^{-1}\to b\in B$, then
\begin{equation}
 g_nx\longrightarrow a
 \quad\text{uniformly for $x$ in each compact subset of }X\setminus\{b\}.
 \label{eq:dyn-convergence}
\end{equation}
Consequently the action of $\Gamma$ on the space of distinct triples
of $B$ is properly discontinuous; in particular, its action on $B$ is
a convergence action.
\end{proposition}

\begin{proof}
Suppose \eqref{eq:dyn-convergence} fails. After taking a subsequence,
there are $x_n\to r\ne b$ with $g_nx_n\to s\ne a$. For each $n$,
density of the vertices and continuity of the particular map $g_n$
give a vertex $v_n$ such that
\[
 \rho(v_n,x_n)<1/n,
 \qquad \rho(g_nv_n,g_nx_n)<1/n.
\]
Thus $v_n\to r$ and $g_nv_n\to s$.

Choose a coordinate separating $r$ from $b$, and another separating
$s$ from $a$. Apply Lemma~\ref{dyn:walls} to these two coordinates.
Choose regular levels strictly between the respective values and
outside its null exceptional sets. Call the resulting cuts the first
and second cuts. The translated first cut has all endpoints uniformly
near $a$, while the inverse-translated second cut has all endpoints
uniformly near $b$.

Let $C_n$ be the graph component of the $r$-side of the first cut
containing $v_n$. Write $F$ for its coordinate and $t$ for its level;
for definiteness suppose $F(b)<t<F(r)$. A neighborhood of $b$ then
misses the entire $r$-side, whose vertices satisfy $F>t$.
Both sides have vertices. A path from $v_n$ to the
other side has a first straddling edge; its endpoint $w_n$ on the
$r$-side lies in $C_n$. Since the inverse-translated second cut is
uniformly near $b$, none of its edges can have an endpoint on the
$r$-side for large $n$. Consequently no edge of $g_nC_n$ straddles the
second cut. Yet $g_nv_n$ lies on its $s$-side, whereas $g_nw_n$ is
uniformly near $a$ and lies on the opposite side. The image of a path
in $C_n$ between $v_n$ and $w_n$ must straddle that cut, a
contradiction.

For proper discontinuity, suppose infinitely many distinct group
elements take a point of one compact set of distinct triples to a
point of another such compact set. Pass to convergent source and
target triples and to subsequences with $g_n\to a$ and
$g_n^{-1}\to b$. The latter limits belong to $B$, since group vertices
are isolated. At least two coordinates of the limiting source triple
differ from $b$. Equation~\eqref{eq:dyn-convergence} sends both to
$a$, contradicting distinctness of the limiting target triple.
\end{proof}

For a convergence action on a compact Hausdorff space with at least
three points, every infinite-order element is either \emph{parabolic},
with one fixed point, or \emph{loxodromic}, with two fixed points
\cite[Lemma~6.2]{Bowditch1999Treelike}. Moreover, a point stabilizer
containing a loxodromic element is virtually cyclic
\cite[Proposition~6.4 and its preceding paragraph]{Bowditch1999Treelike}.
These statements apply because $B$ is a nondegenerate continuum.

\begin{lemma}[Parabolic comparison]\label{dyn:parabolic}
Let $h\in\Gamma$ be parabolic with fixed point $p\in B$. Then
$h^n\to p$ and $h^{-n}\to p$ as vertices of $X$. Both sequences of
maps converge to $p$ uniformly on compact subsets of $X\setminus\{p\}$
and pointwise on all of $B$. Furthermore,
\begin{equation}
 \rho(g,gp)\longrightarrow0\qquad(g\longrightarrow\infty\text{ in }\Gamma).
 \label{eq:dyn-parabolic-comparison}
\end{equation}
\end{lemma}

\begin{proof}
The powers $h^n$ escape every finite vertex set as $n\to+\infty$ or
$n\to-\infty$. If a subsequence tends to $z\in B$,
\eqref{eq:dyn-right-move} gives
$\rho(h^n,h^{n+1})\to0$, whereas continuity of left multiplication by
$h$ gives $h^{n+1}\to hz$. Thus $hz=z$, and uniqueness of the
parabolic fixed point gives $z=p$. Compactness proves convergence of
the entire sequences. Proposition~\ref{dyn:convergence} now supplies
their compact-uniform convergence off $p$; they fix $p$ itself.

If \eqref{eq:dyn-parabolic-comparison} fails, take a subsequence with
$g_n\to a$ and $g_np\to c\ne a$. The cosets $g_n\langle h\rangle$
escape every finite vertex set. Otherwise some fixed $w$ belongs to
infinitely many of them, so $g_n=wh^{j_n}$ on a subsequence, with
$|j_n|\to\infty$. Then both $g_n$ and $g_np$ tend to $wp$, a
contradiction.

For fixed $n$, the sequence $u_{n,j}=g_nh^j$, $j\ge0$, starts at
$g_n$ and tends to $g_np$. By whole-coset escape and
\eqref{eq:dyn-right-move},
\[
 \sup_{j\ge0}\rho(u_{n,j},u_{n,j+1})\longrightarrow0.
\]
Choose three distinct numbers strictly between $0$ and $\rho(a,c)$.
For each number choose the first index $j$ at which
$\rho(a,u_{n,j})$ reaches it. The vanishing step size and compactness
provide, after a common subsequence, three selected vertex sequences
with distinct limits $z_1,z_2,z_3\in B$.

The conjugate $k_n=g_nhg_n^{-1}$ sends each selected vertex to its
successor, so asymptotically fixes these three sequences. If $k_n$
escapes finite sets, pass to limits of $k_n$ and $k_n^{-1}$ in $B$.
At least two of the $z_i$ avoid the repelling limit;
Proposition~\ref{dyn:convergence} would send both to the same
attracting limit, a contradiction. Otherwise take a constant
subsequence of the $k_n$. By continuity, that fixed conjugate of $h$
fixes all three $z_i$, again a contradiction. This proves
\eqref{eq:dyn-parabolic-comparison}.
\end{proof}

\begin{proposition}\label{dyn:minimal}
The action of $\Gamma$ on $B$ is minimal: every orbit is dense.
\end{proposition}

\begin{proof}
Every invariant subset with at least two points is dense. Indeed,
given $a\in B$, choose vertices $g_n\to a$ and pass to
$g_n^{-1}\to b\in B$. An invariant subset with at least two points
contains $z\ne b$, and Proposition~\ref{dyn:convergence} gives
$g_nz\to a$.

It remains to exclude a global fixed point $c$. Choose $a\in B$
different from $c$ and vertices $g_n\to a$. After passing to a
subsequence, Proposition~\ref{dyn:convergence} and $g_nc=c$ force
$g_n^{-1}\to c$. Choose disjoint nonempty closed neighborhoods $A,C$
of $a,c$ in $B$. For large $n$,
\[
 g_nA\subset A,\qquad g_n^{-1}C\subset C.
\]
Fix one such nonidentity $g_n$. Since $\Gamma$ is torsion-free, it has
infinite order. It cannot be parabolic: Lemma~\ref{dyn:parabolic}
would make the forward orbit of a point of $A$ and the backward
orbit of a point of $C$ tend to the same point, which would belong to
both disjoint closed sets. It is therefore loxodromic. Since all of
$\Gamma$ fixes $c$, the stabilizer statement cited above would make
$\Gamma$ virtually cyclic. An infinite virtually cyclic group has
two ends, contrary to the established one-endedness of $\Gamma$.
\end{proof}
